\documentclass[10pt,a4paper]{amsart}
\usepackage[margin=19mm]{geometry}
\usepackage{amsmath,amssymb,mathtools}
\usepackage[scr=rsfs,cal=euler]{mathalfa}
\usepackage{xfrac}
\usepackage[unicode,hidelinks,bookmarks=false]{hyperref}
\newcommand{\ie}{i.e.\ }
\newcommand{\cf}{cf.\ }
\newcommand{\resp}{respectively\ }
\newcommand{\Subsection}{\subsection}
\providecommand{\nobreakdash}{\nobreak\hbox{-}}
\setlength{\emergencystretch}{2em}
\multlinegap0pt
\usepackage{bm}
\usepackage{bbm}
\usepackage{dsfont}
\usepackage{xspace}
\usepackage{cancel}
\usepackage{mathtools}
\usepackage{amsthm}
\makeatletter
\@ifundefined{theorem}{\newtheorem{theorem}{Theorem}}{}
\makeatother
\title[An Iterative Methodology for Unitary Quantum Channel Search]{An Iterative Methodology for Unitary Quantum Channel Search}
\author{Matthew M. Lin, Hao-Wei Huang, Bing-Ze Lu}
\date{Source: arXiv:2506.21455v1 (CC BY 4.0)}
\begin{document}
\maketitle

\noindent\emph{Source and licence.} An Iterative Methodology for Unitary Quantum Channel Search. Version arXiv:2506.21455v1 (CC BY 4.0), distributed under the Creative Commons Attribution 4.0 International licence, as recorded in the arXiv licence metadata for this exact version and shown on the abstract page https://arxiv.org/abs/2506.21455v1. Full rights notice: licenses/arxiv-2506.21455-CC-BY-4.0.txt.\par\smallskip

\noindent\textit{Editorial scope.} Counted unit: the Theorem whose source label is \texttt{thm:global-theorem} in the article (source0.tex lines 376--385, source label \texttt{thm:global-theorem}), delivered together with the article's own complete original proof of it (source0.tex lines 386--424, one \texttt{proof} environment of 39 lines). No mathematics is written, repaired or re-derived: statement and proof are the article's own text line for line. Required context retained uncounted: none. Every definition, display and label of those lines is retained unchanged. Omitted from this package and not redistributed: 1--375, 425--1282. Nothing inside a retained excerpt is omitted or altered: every retained body is delivered exactly as the article's lines are written, so no omission receipt of any kind accompanies this package. Citations: the retained excerpts carry no citation command at all, so no bibliography entry of the article is transcribed and no reference is redistributed. Reference adaptation: none was needed and none was made; every reference inside the delivered unit resolves inside it, so no unresolved reference or citation is produced. Printed slips kept literal: no printed word, symbol, hypothesis or number of the counted statement or of its proof was altered. Layout and class: the article's own class is \texttt{sn-jnl} with packages including epstopdf, graphicx, multirow, amsmath, amssymb, amsfonts, amsthm, mathrsfs, appendix, xcolor, textcomp, manyfoot, booktabs, algorithm; the wrapper is the lane's pinned \texttt{amsart} editorial wrapper at 10pt on A4, which restates the theorem-like environments the retained text needs and adds the article's own macro definitions that the retained text needs (none), with extra wrapper packages (bm, bbm, dsfont, xspace, cancel, mathtools). The delivered numbering is the unit's own. Assets: no retained span contains a figure environment, a graphics-inclusion command, a PDF-page inclusion, a table or a TikZ picture; no image file is copied into this package, the package declares no asset dependency and no image file is redistributed with it. Licence: version arXiv:2506.21455v1 of this article is distributed under the Creative Commons Attribution 4.0 International licence, as recorded in the arXiv licence metadata for this exact version and shown on the abstract page https://arxiv.org/abs/2506.21455v1. A full rights notice is delivered at licenses/arxiv-2506.21455-CC-BY-4.0.txt. Characters: every retained excerpt is pure ASCII, so the delivered engine, XeLaTeX with its default Latin Modern text font, reports no missing character.
    \begin{theorem}\label{thm:global-theorem}
        Suppose $U$ and $V$ are unitary matrices such that
        \begin{eqnarray*}
            U\rho U^* = V\rho V^*
        \end{eqnarray*}
        for every Hermitian matrix $\rho$ with $\mbox{Tr}(\rho)=1$. Then, there exists a scalar $\mu$ with $|\mu|=1$ such that
        \begin{eqnarray*}
            U = \mu V.
        \end{eqnarray*}
    \end{theorem}

    \begin{proof}
        Since
        \begin{eqnarray*}
        U\rho U^* = V\rho V^*
        \end{eqnarray*}
        for all such $\rho$, define the matrix $M = V^* U$, then we have  $M\rho = \rho M$ 
        for every Hermitian matrix $\rho$ with its trace equal to 1.
        
        Choose the standard basis $\{e_i\}_{i=1}^n$ of $\mathbb{R}^n$. Any Hermitian matrix $\rho$ with trace $1$ can be expressed on this basis as
        \begin{eqnarray*}
        \rho = \sum_{i=1}^n \rho_{ii}\, e_i e_i^T + \sum_{i<j} \rho_{ij}\, (e_i e_j^T + e_j e_i^T),
        \end{eqnarray*}
        with $\sum_{i=1}^n \rho_{ii}=1$ and each $\rho_{ij}\in\mathbb{R}$.
        %\\
        Express $M$ in the same basis:
        \begin{eqnarray*}
        M = \sum_{p,q=1}^n M_{pq}\, e_p e_q^T.
        \end{eqnarray*}
        Now, the commutation relation $M\rho = \rho M$ must hold for all such $\rho$. First, we are allowed to consider the effect on the matrices $e_i e_i^T$:
        \begin{eqnarray*}
        M\, e_i e_i^T = e_i e_i^T\, M.
        \end{eqnarray*}
        A straightforward calculation using the above representations shows that this forces $M_{pq} = 0$ whenever $p \neq q$; that is, $M$ must be diagonal.\\
        Next, look at the matrices $e_i e_j^T + e_j e_i^T$ for $i\neq j$. The relation
        \begin{eqnarray*}
        M(e_i e_j^T+e_j e_i^T) = (e_i e_j^T+e_j e_i^T) M
        \end{eqnarray*}
        implies that the diagonal entries of $M$ satisfy $M_{ii} = M_{jj}$ for all $i, j$.\\
        Finally, because $U$ and $V$ are unitary, so is $M = V^* U$. Hence, we have
        \begin{eqnarray*}
        MM^* = I,
        \end{eqnarray*}
        which implies that each diagonal entry of $M$ must have absolute value one; that is, $|M_{ii}|=1$.\\
        Since $M$ is diagonal with all entries equal to the same complex number of modulus one, we can write $M=\mu I$ with $|\mu|=1$. Therefore,
        \begin{eqnarray*}
        V^* U = \mu I \quad \Longrightarrow \quad U = \mu V,
        \end{eqnarray*}
        completing the proof.
    \end{proof}

\end{document}
